\begin{env}[2.1.10]
\phantomsection
\label{II.2.1.10-ko}
$S_+$의 부분집합 $\mathfrak{J}$가 $S$의 아이디얼이면 이를
\emph{$S_+$의 아이디얼}이라 한다. 또한 $\mathfrak{J}$가 $S_+$와,
\emph{$S_+$를 포함하지 않는} $S$의 등급 소아이디얼의 교집합이면 이를
\emph{$S_+$의 등급 소아이디얼}이라 한다(더욱이 이 소아이디얼은
(\hyperref[II.2.1.9-ko]{2.1.9})에 따라 유일하다). $\mathfrak{J}$가
$S_+$의 아이디얼이면, \emph{$S_+$에서의 $\mathfrak{J}$의 근기}는 어떤
거듭제곱이 $\mathfrak{J}$에 속하는 $S_+$의 원소들로 이루어진 집합이다.
다시 말해 이는 $r_+(\mathfrak{J})=r(\mathfrak{J})\cap S_+$이다. 특히
$S_+$에서 $0$의 근기는 다시 $S_+$의 \emph{멱영근기}라 하고
$\mathfrak{N}_+$로 나타낸다. 따라서 이는 $S_+$의 멱영원 전체의
집합이다. $\mathfrak{J}$가 $S_+$의 \emph{등급} 아이디얼이면 그 근기
$r_+(\mathfrak{J})$도 등급 아이디얼이다. 실제로 몫환
$S/\mathfrak{J}$로 넘어가면 $\mathfrak{J}=0$인 경우로 환원할 수 있고,
$x=x_h+x_{h+1}+\cdots+x_k$가 멱영원이면
$x_i\in S_i$ ($1\leq h\leq i\leq k$)들도 모두 멱영원임을 보이면
충분하다. $x_k\neq 0$이라고 가정할 수 있고, 이때 $x^n$의 최고차
성분은 $x_k^n$이므로 $x_k$는 멱영원이다. 이제 $k$에 대한 귀납법을
적용한다. 등급환 $S$가 \emph{본질적으로 축소}라는 것은
$\mathfrak{N}_+=0$이라는 뜻이다. 다시 말해 $S_+$가 $\neq 0$인 멱영원을
포함하지 않는다는 뜻이다.
\end{env}

\begin{env}[2.1.11]
\phantomsection
\label{II.2.1.11-ko}
등급환 $S$에서 원소 $x$가 어떤 영이 아닌 원소와 곱하여 $0$이 되는
영인자이면 그 최고차 동차성분도 영인자임에 유의하자. 환 $S$가
\emph{본질적으로 정역}이라는 것은 $S_+$가
(\emph{단위원을 갖지 않는 환으로서}) 영인자를 포함하지 않아 영이 아닌
두 원소의 곱이 $0$일 수 없고, 또한 $\neq 0$인 것을 뜻한다. 이를 위해서는
$S_+$의 동차 원소 가운데 $\neq 0$인 것은 모두 이 환에서 영인자가 아니어서
어떤 영이 아닌 원소와 곱해도 $0$이 되지 않으면 충분하다.
$\mathfrak{p}$가 $S_+$의
등급 소아이디얼이면 $S/\mathfrak{p}$가 본질적으로 정역임은 명백하다.

$S$를 본질적으로 정역인 등급환이라 하고 $x_0\in S_0$라 하자. 이때
$S_+$의 동차 원소 $f\neq 0$가 \emph{하나라도} 존재하여
$x_0f=0$이면 $x_0S_+=0$이다. 실제로 모든 $g\in S_+$에
대하여 $(x_0g)f=(x_0f)g=0$이고, 가정에 따라 $x_0g=0$이다. 따라서
$S$가 정역이기 위한 필요충분조건은 $S_0$가 정역이고 $S_0$에서
$S_+$의 소멸자가 $0$뿐인 것이다.
\end{env}

\subsection{등급환의 분수환.}
\label{subsection:II.2.2-ko}

\begin{env}[2.2.1]
\phantomsection
\label{II.2.2.1-ko}
$S$를 음이 아닌 차수로 등급이 주어진 환이라 하고, $f$를 차수가
$d>0$인 $S$의 \emph{동차} 원소라 하자. 그러면 분수환 $S'=S_f$는
$S'_n$을 $x\in S_{n+kd}$이고 $k\geq 0$일 때의 $x/f^k$ 전체의
집합으로 놓음으로써 등급환이 된다(여기서 $n$은 임의의 음의 값도 취할
수 있음에 유의하자). $S'$의 \emph{차수가 $0$인} 원소들로 이루어진
부분환 $S'_0=(S_f)_0$을 $S_{(f)}$로 나타내겠다.

$f\in S_d$이면 $S_f$의 단항식 $(f/1)^h$들($h$는 영을 포함한 임의의
정수)은 환 $S_{(f)}$ 위의 \emph{자유계}를 이루며, 그 선형결합 전체의
집합은 바로
\oldpage[II]{24}
환 $(S^{(d)})_f$이다. 따라서 이 환은
\emph{$S_{(f)}[T,T^{-1}]=S_{(f)}\otimes_{\mathbf{Z}}\mathbf{Z}[T,T^{-1}]$와
동형이다}($T$는 부정원이다). 실제로
$\sum_{h=-a}^{b}z_h(f/1)^h=0$인 관계가 있고 $z_h=x_h/f^m$이며 모든
$x_h$가 $S_{md}$에 속한다고 하자. 정의에 따라 이 관계는
어떤 $k\geq a$가 존재하여
$\sum_{h=-a}^{b}f^{h+k}x_h=0$이 되는 것과 동치이다. 이 합의 항들은
차수가 서로 다르므로 모든 $h$에 대하여 $f^{h+k}x_h=0$이고, 따라서
모든 $h$에 대하여 $z_h=0$이다.

$M$이 등급 $S$-가군이면 $M'=M_f$는 등급 $S_f$-가군이다. 여기서
$M'_n$은 $z/f^k$ 꼴의 원소들로 이루어진 집합이며, 여기서
$z\in M_{n+kd}$이고 $k\geq 0$이다.
$M'$의 차수가 $0$인 동차 원소들의 집합을 $M_{(f)}$로 나타낸다.
$M_{(f)}$가 $S_{(f)}$-가군이고
$(M^{(d)})_f=M_{(f)}\otimes_{S_{(f)}}(S^{(d)})_f$임은 바로 알 수 있다.
\end{env}
